\chapter{Implementation and Testing}

This chapter documents the technical implementation details, algorithms, external dependencies, and testing strategies employed in the MyAIStorybook system. The implementation demonstrates practical application of AI technologies and software engineering best practices to achieve reliable storybook generation.

\section{Algorithm Design}

The system employs several core algorithms to orchestrate story generation, ensure quality, and maintain consistency across generated content.

\subsection{Multi-Agent Story Generation Pipeline}

The multi-agent pipeline coordinates sequential processing through specialized agents to produce high-quality stories.

\begin{verbatim}
ALGORITHM: Multi-Agent Story Generation
INPUT: user_prompt (string), generate_images (boolean)
OUTPUT: complete_story (JSON), status (string)

BEGIN
  1. Initialize timestamp = current_time()
  2. Call PromptAgent.process_prompt(user_prompt)
     - Validate prompt length (5-500 chars)
     - Classify as: normal, detailed, invalid, nonsense
     - Enhance simple prompts with context
     - RETURN enhanced_prompt, prompt_type
  
  3. IF prompt_type IN {invalid, nonsense} THEN
       RETURN error_response("Invalid prompt")
  
  4. Call DirectorAgent.create_story(enhanced_prompt)
     4.1. WriterAgent.generate_story(enhanced_prompt)
          - Format LLM prompt with story structure template
          - Call Ollama LLM (llama3.1:8b-instruct-q4_K_M)
          - Parse JSON response (retry up to 2 times if fails)
          - Extract: title, setting, characters, scenes
          - RETURN story_draft
     
     4.2. ReviewerAgent.review_story(story_draft)
          - Validate coherence and consistency
          - Check age-appropriateness
          - Assess educational value
          - RETURN reviewed_story
     
     4.3. EditorAgent.edit_story(reviewed_story)
          - Refine language and formatting
          - Generate PDF using ReportLab
          - Save to exports/ directory
          - RETURN final_story
  
  5. IF generate_images = TRUE THEN
     5.1. Pause Ollama service (free GPU memory)
     5.2. FOR each scene IN final_story.scenes DO
            - Truncate prompt to 77 tokens
            - Call ImageAgent.generate_image(scene.description)
            - Save PNG to images/ directory
            - Update scene.image_path
     5.3. Resume Ollama service
  
  6. Save final_story to stories/ directory
  7. Trigger EvaluationAgent in background (async)
  8. RETURN final_story, status
END
\end{verbatim}

\subsection{Character Consistency Algorithm}

The IP-Adapter integration ensures consistent character appearance across all story illustrations.

\begin{verbatim}
ALGORITHM: Character Consistency with IP-Adapter
INPUT: user_photo (image), scene_descriptions (list)
OUTPUT: consistent_images (list of image paths)

BEGIN
  1. Load user_photo and preprocess
     - Resize to 512x512
     - Convert to RGB
     - Normalize pixel values
  
  2. Initialize PersonalizedImageAgent with user_photo
     - Load IP-Adapter weights
     - Encode user_photo to embedding
     - Store reference embedding
  
  3. FOR each scene_description IN scene_descriptions DO
     3.1. Construct enhanced_prompt:
          - Combine scene_description with style keywords
          - Add character consistency tags
          - Apply negative prompts (distortion, multiple faces)
     
     3.2. First pass: txt2img generation
          - Use DPMSolver scheduler (Karras sigmas)
          - Resolution: 512x512
          - Inference steps: 25
          - Guidance scale: 7.5
          - Inject IP-Adapter embedding
     
     3.3. Second pass: img2img refinement
          - Input: first_pass_image
          - Denoise strength: 0.3
          - Resolution: 768x768 (upscale)
          - Maintain character features
     
     3.4. Save final_image with timestamp
     3.5. ADD final_image_path to consistent_images
  
  4. RETURN consistent_images
END
\end{verbatim}

\subsection{Conversation Context Management}

The workshop agent maintains conversation context for iterative story development.

\begin{verbatim}
ALGORITHM: Workshop Conversation Management
INPUT: session_id, user_message
OUTPUT: agent_response, ready_to_generate

BEGIN
  1. Load session from database by session_id
  2. Load conversation_history (ordered by timestamp)
  3. Extract metadata from previous messages:
     - story_requirements
     - collected_information
     - conversation_state
  
  4. Analyze user_message:
     - Extract story elements (characters, setting, theme)
     - Identify requirements satisfied
     - Update metadata
  
  5. Determine conversation phase:
     IF mode = "new_idea" THEN
        - Check if essential fields complete
        - ready = (theme AND characters AND setting)
     ELSE IF mode = "improvement" THEN
        - Check if user story provided
        - ready = (user_story.length > 0)
  
  6. Generate contextual response:
     - Format prompt with conversation history
     - Call LLM with context window
     - Generate follow-up questions OR confirmation
  
  7. Save user_message and agent_response to database
  8. RETURN agent_response, ready_to_generate
END
\end{verbatim}

\section{External APIs/SDKs}

The MyAIStorybook system integrates several external APIs, SDKs, and libraries to provide its comprehensive functionality.

\begin{table}[H]
\centering
\caption{External APIs/SDKs Used in Implementation}
\begin{tabular}{|p{3.5cm}|p{4.5cm}|p{3cm}|p{3cm}|}
\hline
\textbf{API/SDK and Version} & \textbf{Description} & \textbf{Purpose of Usage} & \textbf{Key Endpoints/Classes} \\ \hline
Ollama (Latest) & Local LLM server for running quantized language models & Story text generation, chat responses & ollama run llama3.1:8b-instruct-q4\_K\_M \\ \hline
stable-diffusion-webui & AUTOMATIC1111's web UI with local API & IP-Adapter integration, personalized image generation & /sdapi/v1/txt2img, /sdapi/v1/img2img \\ \hline
FastAPI 0.100+ & Modern async web framework for Python & Backend REST API server & @app.post(), @app.get() decorators \\ \hline
Diffusers 0.21+ & Hugging Face library for diffusion models & Fallback image generation pipeline & StableDiffusionPipeline, DPMSolverMultistepScheduler \\ \hline
PyTorch 2.0+ & Deep learning framework & Model inference, tensor operations & torch.cuda, torch.load() \\ \hline
SQLAlchemy 2.0+ & Python SQL toolkit and ORM & Database operations & Session, Base.metadata.create\_all() \\ \hline
ReportLab 4.0+ & PDF generation library & Story PDF export & canvas.Canvas, SimpleDocTemplate \\ \hline
Next.js 14+ & React framework & Frontend application & pages/, app/ routing \\ \hline
PostgreSQL & Relational database & User auth, story metadata & N/A (via SQLAlchemy) \\ \hline
IP-Adapter & Image prompt adapter for SD & Character consistency & Via WebUI API integration \\ \hline
\end{tabular}
\label{tab:apis_sdks}
\end{table}

\section{Testing Details}

The system employs a combination of unit testing for individual components and integration testing for the complete pipeline. Testing focuses on ensuring reliability, correctness, and graceful error handling.

\subsection{Unit Testing}

Each agent module is tested independently to verify correct functionality and error handling.

\textbf{PromptAgent Testing:}

\begin{verbatim}
Test Case: test_prompt_validation()
Purpose: Verify prompt classification and enhancement
Input: Various prompts (valid, short, long, nonsense)
Expected Output: Correct classification and enhancement
Result: PASS - All prompt types correctly identified

Test Case: test_prompt_length_validation()
Purpose: Ensure prompts within 5-500 character range accepted
Input: Prompts of varying lengths
Expected Output: Accept valid lengths, reject invalid
Result: PASS - Boundary conditions handled correctly
\end{verbatim}

\textbf{WriterAgent Testing:}

\begin{verbatim}
Test Case: test_story_generation()
Purpose: Verify story JSON structure and completeness
Input: Enhanced prompt "A brave knight's adventure"
Expected Output: Valid JSON with title, setting, characters, 
                 scenes
Result: PASS - Generated story contains all required fields

Test Case: test_json_retry_mechanism()
Purpose: Verify retry logic on malformed JSON
Input: Prompt causing occasional JSON errors
Expected Output: Successful generation after retry
Result: PASS - Retry mechanism recovers from parsing errors
\end{verbatim}

\textbf{ImageAgent Testing:}

\begin{verbatim}
Test Case: test_image_generation_gpu()
Purpose: Verify image generation on GPU-enabled systems
Input: Scene description "A castle on a hill"
Expected Output: PNG image saved to images/ directory
Result: PASS - Image generated in ~40 seconds (GTX 1060)

Test Case: test_image_generation_cpu_fallback()
Purpose: Verify CPU fallback when GPU unavailable
Input: Same scene description
Expected Output: Image generated with CPU parameters
Result: PASS - Generation completes in ~2.5 minutes
\end{verbatim}

\subsection{Integration Testing}

Integration tests verify the complete multi-agent pipeline and frontend-backend communication.

\textbf{End-to-End Story Generation:}

\begin{verbatim}
Test Case: test_complete_generation_pipeline()
Purpose: Verify full story generation from prompt to PDF
Input: User prompt via /api/generate endpoint
Steps:
  1. POST request to /api/generate with test prompt
  2. Verify response contains story JSON
  3. Check all scenes have text content
  4. Verify images generated if enabled
  5. Confirm PDF created in exports/
Expected Output: Complete story with all components
Result: PASS - Pipeline executes successfully
Average Time: 45 seconds (with images)
\end{verbatim}

\textbf{Error Handling Testing:}

\begin{verbatim}
Test Case: test_graceful_degradation()
Purpose: Verify system continues with text-only on image failure
Input: Valid prompt with image generation enabled
Simulation: Force image generation error
Expected Output: Story generated without images, no crash
Result: PASS - System delivers text-only story with warning

Test Case: test_invalid_prompt_handling()
Purpose: Verify appropriate error response for invalid input
Input: Nonsense prompt "asdfghjkl"
Expected Output: HTTP 400 with error message
Result: PASS - Clear error message returned to user
\end{verbatim}

\subsection{Performance Testing}

Performance tests measure system response times under typical usage conditions.

\begin{verbatim}
Test Results (NVIDIA GTX 1060, Intel i5, 16GB RAM):
- Story text generation: 12-15 seconds
- Single image generation (GPU): 35-45 seconds
- Complete 3-scene story with images: 2-2.5 minutes
- PDF export: 2-3 seconds
- Database query latency: <100ms

Test Results (CPU-only, Intel i7, 16GB RAM):
- Story text generation: 18-22 seconds
- Single image generation (CPU): 2-3 minutes
- Complete 3-scene story with images: 8-10 minutes
\end{verbatim}
